\section{Methods}
\label{sec:methods}

% This is the section that earns or loses the paper. Everything here should be
% traceable to a workflow step (WF-xxx) in the planning workbook, and every
% subsection should end able to answer "how would someone reproduce this?"

% ---------------------------------------------------------------------------
% AS-BUILT OUTLINES (added 2026-09-02). Each \OUTLINE block below records what
% was actually run, in note form, so the prose can be written from it. Notes are
% raw material, not draft sentences. Delete each block once its subsection is
% written.
% ---------------------------------------------------------------------------

\subsection{Study area and analysis extent}
\label{sec:extent}

The study area was defined as the southern African cheetah range plus a buffered analysis extent. All layers were aligned to a fixed equal-area projection and a \cellsize\ snap grid, so every raster is directly comparable cell for cell. Cores were derived from the historical confirmed-range baseline and held constant across 2012, 2016, 2020, and 2024. The final network consisted of 23 cores and 45 selected neighboring core pairs.

\subsection{Rationale for the working resolution}
\label{sec:scale}

I used a 1-km working resolution because it represents a compromise between the coarse 10-km occurrence baseline and the finer environmental covariates, which range from 30 m to 1 km. This resolution is fine enough to identify meaningful pinch points in the landscape while avoiding the implication of precision that the occurrence evidence cannot support.

\subsection{Range cores}
\label{sec:cores}
% Sentence starters: "I defined a core area as ..." / "Cells were grouped into patches using ..." / "I retained only patches larger than ..."
\TODO{WF-011. Three or more core definitions from the archived confirmed-range data,
differing in evidence-category threshold and minimum patch area. Report how many
cores each definition yields. State explicitly that cores are HELD CONSTANT across
all four temporal snapshots, and why (CAU-002): no comparable range reassessment
exists after the baseline period, so any core movement would be invented.}

I defined a core area as a contiguous patch of cells that were confirmed presence in the historical baseline. Cells were grouped into patches using an 8-neighbor rule, and only patches larger than a minimum area threshold were retained. I used three core definitions from the archived confirmed-range data, differing in evidence-category threshold and minimum patch area. The final network consisted of 23 cores, which were held constant across all four temporal snapshots to avoid inventing core movement in the absence of comparable range reassessment.

\begin{outline}
I represented the 23 baseline-defined core areas as fixed nodes across all four temporal snapshots. Candidate links were defined as the union of an undirected three-nearest-neighbor graph based on Euclidean distances between core centroids and the graph’s minimum spanning tree. Removing duplicate undirected pairs yielded 45 selected core pairs. These straight-line links defined the pairs evaluated using least-cost analysis and were not interpreted as observed movement corridors.
\end{outline}

\subsection{Covariates and their temporal treatment}
\label{sec:covariates}
% Sentence starters: "The predictor layers represented ..." / "I treated a variable as dynamic when ..." / "I treated a variable as static when ..."
\TODO{Table~\ref{tab:covariates} lists every layer, its native resolution, its
source year(s), and critically whether it is treated as DYNAMIC (varies across
snapshots) or STATIC (held fixed, contributing no temporal signal). Do not bury this
distinction. Roads, livestock, terrain and hydrography are static; built-up fraction,
nighttime radiance and continuous vegetation fields are dynamic.}

\begin{outline}
\begin{itemize}
  \item Built-up fraction: GHSL, 2012 / 2016 / 2020 / 2024.
  \item Nighttime lights: VIIRS, flare-masked before differencing.
  \item Land cover: MODIS MCD12Q1, reclassified to open vegetation, woody, cropland,
        built, water, bare.
  \item Vegetation continuous fields (VCF): tree cover, non-tree vegetation, bare ground.
  \item Roads: GRIP4 Africa, clipped to the study extent.
  \item Livestock: GLW4 cattle, goats and sheep, combined into a livestock-exposure index.
  \item Terrain: approximately 30\,m SRTM DEM pulled through Google Earth Engine and
        aggregated to the \cellsize\ grid. Derivatives: mean slope, maximum slope,
        terrain resistance surface.
  \item Veterinary fences: WWF/KAZA reference fence dataset, with main, near-barrier and
        conservative fence scenarios. Kruger fence references treated as scenario
        information, not as universal barriers.
  \item Dynamic: built-up fraction, nighttime lights, tree cover, non-tree vegetation,
        bare-ground VCF.
  \item Static: roads, livestock, terrain, hydrography, fences, protected areas, core
        definitions.
  \item Categorical land cover used for context and masking only, never as an independent
        temporal resistance penalty.
\end{itemize}
\end{outline}

\paragraph{Why categorical land cover is not the change carrier.}
\TODO{CAU-008 and the MODIS class-instability problem. Annual categorical products
flip classes between years for reasons unrelated to landscape change. Temporal signal
is therefore carried on continuous fields; categorical land cover is used for masking
and description only.}

\subsection{Resistance scenarios}
\label{sec:resistance}
% Sentence starters: "Resistance values were assigned to represent ..." / "The terrain component was derived from ..." / "I compared alternative vegetation weights because ..."
\TODO{WF-013, WF-014. Present RES-01 to RES-05 as a table of variable-to-resistance
functions with the evidence or rationale for each. State that the scenario set was
fixed before results were viewed and that no single scenario is designated as
preferred. Report the transformation (linear vs exponential) per scenario and cite
\textcite{zeller2012resistance} for why the transformation matters as much as the
weights.}

\begin{outline}

Terrain component. The DEM was not used directly as resistance; slope was derived from
it. Mean slope carries the primary terrain surface, maximum slope retained for
sensitivity and context. Terrain resistance function: slope $\le 10^\circ$ low
resistance; $10$ to $30^\circ$ linearly increasing; $\ge 30^\circ$ capped high;
resistance rescaled to 1 to 10.

Original primary temporal surface. Anthropogenic component 80\% of total, terrain 20\%.
Within the anthropogenic component: built-up 30\%, roads 30\%, livestock 30\%,
nighttime lights 10\%.

Three internal weighting models compared: primary, equal, infrastructure-emphasis.

Vegetation correction. The first temporal model omitted VCF from the resistance
calculation, which is a mismatch with the predeclared protocol. Three further
sensitivity scenarios were added, keeping the original anthropogenic-only surfaces
intact rather than replacing them:
\begin{itemize}
  \item vegetation-low: 10\% vegetation, 70\% anthropogenic, 20\% terrain
  \item vegetation-balanced: 20\% vegetation, 60\% anthropogenic, 20\% terrain
  \item vegetation-high: 40\% vegetation, 40\% anthropogenic, 20\% terrain
\end{itemize}
The vegetation proxy uses continuous VCF fields to represent sparse and bare cover. It
does not claim to measure fine-scale habitat edge or observed movement preference, and
should be described that way.

Reconciliation still needed: these six weighting models are not the same object as the
predeclared RES-01 to RES-06 scenario set. Either map each one onto a predeclared
scenario ID or state plainly in the text that the implemented set differs and why.
\end{outline}

\subsection{Connectivity modelling}
\label{sec:connectivity}
% Sentence starters: "For each selected core pair, we calculated ..." / "The straight-line link was used to define ..." / "The least-cost path could depart from that line when ..."
\TODO{WF-015, WF-016. Least-cost / effective-distance work in ArcGIS Pro (Distance
Accumulation, Optimal Region Connections, Corridor). Current flow computed with
Circuitscape.jl \parencite{anantharaman2020circuitscape} in pairwise mode; give the
solver, the source/ground configuration, and how current was normalised for
comparability across years. State why the omnidirectional approach
\parencite{landau2021omniscape} was NOT used: cores are defined a priori, so a
core-to-core formulation is the correct one.}

\begin{outline}
\begin{itemize}
  \item Least-cost paths computed between all 45 selected core pairs.
  \item Cost-distance analysis run on a common snap raster, cell size, mask, extent and
        projection, so surfaces are comparable across years and scenarios.
  \item Implementation shortcut worth documenting: one \texttt{CostDistance} run per
        source core, then multiple destination paths extracted from that single
        accumulated surface with \texttt{CostPathAsPolyline}. Equivalent result to
        pairwise runs, avoids recomputing the same source surface 45 times. Say it is
        equivalent and say why.
  \item Circuitscape current-flow stage not yet run. Either report it once complete or
        change this subsection to describe a least-cost-only design.
\end{itemize}
\end{outline}

\paragraph{Corridor definition.}
\TODO{Decide and state ONCE whether corridors are (a) continuous current surfaces
summarised by quantile or (b) thresholded polygons. This choice sets the denominator
for the protected-area comparison in \S\ref{sec:pa} and cannot be made after seeing
the overlap result.}

\subsection{Network structure}
\label{sec:graph}
\TODO{WF-017. Cores as nodes, effective distances as edge weights. Report the
probability of connectivity and per-node $dPC$ \parencite{saura2007pc}, computed in
Conefor \parencite{saura2009conefor} or an equivalent, and betweenness on the same
graph. Justify the metric set against \textcite{keeley2021metrics}. Report
sensitivity to the dispersal-distance parameter; do not report a single value.}

\begin{outline}
edge betweenness computed on the 23-core network, link
importance ranked on normalised edge betweenness. PC and $dPC$ not yet computed, and
dispersal distance is still the 100\,km placeholder.
\end{outline}

\subsection{Temporal comparison}
\label{sec:temporal}
% Sentence starters: "I compared the same core pairs across ..." / "Absolute change was calculated as ..." / "Proportional change was calculated as ..."
\TODO{WF-018. Same cores, same scenarios, same resolution across all four snapshots,
so that differences are attributable to covariate change alone. Give the change
equations. Report both absolute and proportional change.}

\begin{outline}
Completed snapshot reports exist for 2012, 2016, 2020 and
2024. First-pass anthropogenic-only comparison:
\begin{itemize}
  \item all 45 links present in every year
  \item all 32 robust links have both 2012 and 2024 costs
  \item median robust-link change approximately $-0.002\%$
  \item largest robust-link increase approximately $0.349\%$
  \item no robust link changed by $\ge 1\%$
\end{itemize}
This is not the final temporal result, because that model omitted VCF. The
vegetation-weight sensitivity run is in progress and supersedes it. Write the
subsection so the near-zero change finding is framed as a real result about modelled
permeability rather than a null to be explained away, and note that it holds under a
model with no vegetation term.
\end{outline}

\subsection{Protected-area coverage}
\label{sec:pa}
\TODO{WF-019, CAU-011, CAU-012. Name the \wdpaversion. State polygon-only handling
and what was excluded. State the null explicitly: high-current area is compared
against \TODO{the full extent OR the modelled network} using an area-standardised
permutation test. Predeclare this.}

\subsection{Robustness and uncertainty}
\label{sec:uncertainty}
% Sentence starters: "I considered a result robust when ..." / "A link was classified as sensitive when ..." / "Locations with low agreement were treated as ..."
\TODO{WF-022. The full sensitivity grid: \nscen\ resistance scenarios $\times$ core
definitions $\times$ resolutions $\times$ snapshots. Report the resulting run count.
Define "robust" numerically (e.g. a cell in the top decile of current under at least
$k$ of $n$ runs) and fix $k$ in advance. Produce an agreement map. Low-agreement
locations are reported as uncertain, never as priorities.}

\begin{outline}
\begin{itemize}
  \item 90 weight comparisons evaluated.
  \item 34 of 45 core pairs robust under the tested alternative weights, 11 sensitive.
  \item Fence scenarios changed 5 links: 1--13, 1--17, 9--13, 13--17, 27--30.
  \item 32 links classified primary-priority eligible; 13 retained as uncertainty and
        data-collection priorities.
\end{itemize}
Two things to resolve before writing: state the numeric robustness rule actually applied
and whether it matches the predeclared $k/n = 0.8$; and reconcile the 34 robust pairs
with the 32 primary-priority-eligible links, since those are different counts and a
reader will notice.
\end{outline}

\subsection{Independent plausibility assessment}
\label{sec:validation}
\TODO{WF-020, WF-021, CAU-004, CAU-017. See the validation decision note before
writing this. If the plausibility-check route is taken, this subsection is titled as
above and must NOT use the word "validation". If the empirical-resistance route is
taken, add RES-06, describe the fitting procedure, and report country-held-out
performance following \textcite{roberts2017cv}.}

\subsection{Monitoring prioritisation}
\label{sec:mcda}
\TODO{WF-023, WF-024, CAU-020. Components, normalisation, and the predeclared weight
set, plus at least two alternative weightings. Report rank stability (Spearman /
Kendall) across weightings. Publish component scores alongside the composite so a
reader can reweight.}

\begin{outline}
link ranking currently rests on normalised edge
betweenness plus the robust / sensitive split, with 32 primary-priority-eligible and 13
uncertainty-priority links. The full MCDA with three weight sets is not yet built.
\end{outline}

\subsection{Software and reproducibility}
\label{sec:software}
\TODO{ArcGIS Pro \TODO{version}; Circuitscape.jl \TODO{version}; Julia \TODO{version};
R \TODO{version} or Python \TODO{version}; Google Earth Engine for acquisition and
harmonisation. All geoprocessing scripted in arcpy rather than ModelBuilder. Archive
location and what can and cannot be redistributed (CAU-012, CAU-013).}

\begin{outline}
ArcGIS Pro Spatial Analyst via arcpy
(\texttt{CostDistance}, \texttt{CostPathAsPolyline}); Google Earth Engine for SRTM,
GHSL, VIIRS, MODIS MCD12Q1 and VCF acquisition and harmonisation. Record versions.
\end{outline}
